% !TEX root = ../main.tex
% Бүлэг 2

\chapter{Холбоотой бүтээлүүдийн судалгаа}
\label{Chapter2}
\pagecolor{Beige}

%==============================================================
\section{Ясны хугарал ба рентген оношилгооны үндэс}
%==============================================================

\subsection*{Ясны хугарлын ангилал}

Ясны хугарал (bone fracture) нь ясны бүтцийн бүрэн бүтэн байдал алдагдах эмгэг өөрчлөлт бөгөөд гадны үйлчилж буй хүч нь ясны механик тэсвэрлэх чадвараас давсан үед үүсдэг~\cite{ref22}. Хугарлыг арьсны гэмтлийн байдлаар нь нээлттэй (open/compound) болон хаалттай (closed/simple) гэж ангилдаг. Нээлттэй хугарлын үед яс арьсыг нэвтлэн гарсан байдаг бол хаалттай хугарлын үед арьсны бүтэн байдал хадгалагдсан байдаг~\cite{ref22}.

Хугарлын шугамын хэлбэр, чиглэлээс хамааран хөндлөн (transverse), ташуу (oblique), мушгирсан (spiral), олон хэлтэрхийт (comminuted), дарагдсан (compression), ногоон мөчир хэлбэрийн (greenstick) болон стрессийн (stress) хугарал зэрэг үндсэн төрлүүдэд хуваадаг. Эдгээр төрлүүд нь үүсэх механизм, гэмтлийн хүчний чиглэл болон ясанд үүссэн өөрчлөлтөөрөө ялгаатай байдаг~\cite{ref22,ref23}.

Мөн хугарлыг үүсэх шалтгаанаар нь гэмтлийн (traumatic), давтагдах ачааллын (stress) болон эмгэг (pathological) хугарал гэж ангилна. Эмгэг хугарал нь остеопороз, хавдар зэрэг ясны бат бөх чанарыг бууруулдаг өвчний улмаас үүсдэг онцлогтой~\cite{ref24}.

Олон улсын хэмжээнд өргөн хэрэглэгддэг OTA/AO ангиллын систем нь хугарлыг ясны төрөл, анатомийн байршил болон хугарлын хэв шинжээр нь ангилдаг. Энэхүү ангилал нь хугарлын оношилгоо, эмчилгээний төлөвлөлт болон судалгааны ажлуудад түгээмэл хэрэглэгддэг стандарт юм~\cite{ref25}.

\subsection*{Рентген зургийн оношилгооны үндэс}

Рентген зураг (radiograph) нь ясны хугарлыг илрүүлэхэд хамгийн өргөн хэрэглэгддэг анхан шатны дүрс оношилгооны арга юм. Энэхүү арга нь цахилгаан соронзон цацрагийг ашиглан биеийн доторх яс, эд эрхтний бүтцийг дүрслэн харуулах боломж олгодог бөгөөд хугарлын байршил, хэлбэр, шилжилтийг үнэлэхэд чухал ач холбогдолтой~\cite{ref09}. Оношилгооны нарийвчлалыг нэмэгдүүлэх зорилгоор ихэвчлэн хоёр болон түүнээс дээш өнцгөөс зураг авдаг.

Рентген шинжилгээ нь хурдан хугацаанд хийгддэг, өртөг багатай, өргөн хүртээмжтэй зэрэг давуу талтай боловч зарим төрлийн хугарлыг илрүүлэхэд хязгаарлалттай байдаг. Ялангуяа нарийн ан цавтай, шилжилт багатай эсвэл эхэн үедээ тод харагдахгүй хугарлууд рентген зураг дээр ялгарахгүй байх тохиолдол түгээмэл байдаг~\cite{ref10}. Судалгаагаар хугарлын оношилгооны алдааны ихэнх хувийг дүрслэл дээр бүдэг харагдах эсвэл огт ялгарахгүй далд хугарал эзэлдэг болохыг тогтоосон байна~\cite{ref26}.

Иймээс рентген зурагт суурилсан оношилгооны нарийвчлалыг нэмэгдүүлэх, эмчийн ажлыг дэмжих зорилгоор компьютерийн хараа болон хиймэл оюунд суурилсан аргуудыг ашиглах судалгаа эрчимтэй хөгжиж байна.

\subsection*{Ясны эдгэрэлтийн үйл явц}

Ясны хугарлын эдгэрэлт нь эсийн болон молекулын түвшинд явагддаг нарийн зохицуулалттай биологийн процесс бөгөөд ерөнхийдөө үрэвслийн үе шат, зөөлөн каллусын үе шат, хатуу каллусын үе шат болон дахин хэлбэршлийн үе шат гэсэн дөрвөн үндсэн үе шатанд хуваагддаг~\cite{ref27}.

Эхний шатанд хугарлын хэсэгт цус алдалт болон үрэвслийн урвал үүсч, эдгэрэлтийн процесс эхэлдэг. Дараагийн шатанд зөөлөн каллус үүсэж хугарсан хэсгийг түр хугацаанд тогтворжуулна. Үүний дараа эрдэсжилтийн үйл явц явагдаж хатуу каллус үүсэх бөгөөд хугарсан ясны механик бат бөх чанар нэмэгддэг. Эцэст нь остеобласт болон остеокласт эсүүдийн үйл ажиллагааны үр дүнд шинэ үүссэн яс аажмаар анхны бүтэц, хэлбэртээ ойртон өөрчлөгдөнө~\cite{ref27,ref28}.

Ясны эдгэрэлтийн хурд, чанарт нас, хоол тэжээлийн байдал, цусны хангамж, механик тогтвортой байдал, тамхидалт болон чихрийн шижин зэрэг хүчин зүйлс чухал нөлөө үзүүлдэг~\cite{ref29}. Ихэнх хугарал 6--12 долоо хоногийн хугацаанд үндсэндээ эдгэрдэг боловч бүрэн дахин хэлбэршил болон үйл ажиллагааны сэргэлт нь хэдэн сар, зарим тохиолдолд хэдэн жил үргэлжилж болно.

Ясны хугарал, түүний дүрс оношилгоо болон эдгэрэлтийн биологийн үндсийг ойлгох нь хугарлыг ангилах, сегментлэх, хугарлын зэргийг тодорхойлох болон эдгэрэлтийн явцыг үнэлэхэд ашиглагдах хиймэл оюуны аргуудын онолын суурийг бүрдүүлдэг.

%==============================================================
\section{Гүн сургалтын онолын үндэс}
%==============================================================

Гүн сургалт (Deep Learning) нь хиймэл нейрон сүлжээнд суурилсан машин сургалтын дэвшилтэт чиглэл бөгөөд их хэмжээний өгөгдлөөс олон түвшний төлөөлөл (representation) автоматаар сурч чаддагаараа онцлогтой~\cite{ref31}. Сүүлийн жилүүдэд гүн сургалтын аргууд нь компьютер хараа, хэл боловсруулах болон эмнэлгийн дүрс боловсруулалтын салбарт өндөр үр дүн үзүүлж байна. Ялангуяа рентген, компьютер томограф (CT), соронзон резонансын дүрслэл (MRI) зэрэг эмнэлгийн дүрсүүдийг боловсруулахад хүний гараар онцлог шинж тодорхойлох уламжлалт аргуудаас илүү үр ашигтай болох нь олон судалгаагаар нотлогдсон.

Энэхүү судалгаанд ясны хугарлыг ангилах, сегментлэх, хугарлын зэргийг үнэлэх болон үр дүнг тайлбарлахад ашиглагдсан үндсэн гүн сургалтын аргуудын онолын үндсийг дараах байдлаар авч үзнэ.

\subsection*{Конволюцийн нейрон сүлжээ}

\begin{figure}[H]
    \centering
    \figauto[0.64\textwidth]{Figures/CNN_LAST_2.1.png}{CNN архитектурын ерөнхий бүтэц (албан ёсны эх сурвалжаас)}
    \figcaption{CNN архитектурын ерөнхий бүтэц}{CNN архитектурын ерөнхий бүтэц — convolution болон pooling давхаргуудаар дүрсээс орон зайн онцлог шинжүүдийг шатлан ялгаж, бүрэн холболтот давхаргаар ангиллын үр дүн гаргадаг}
    \label{fig:cnn_arch}
\end{figure}

Конволюцийн нейрон сүлжээ (Convolutional Neural Network, CNN) нь дүрс боловсруулалтын хамгийн өргөн хэрэглэгддэг архитектурын нэг юм~\cite{lecun1998}. CNN нь дүрсээс ирмэг, бүтэц, хэлбэр зэрэг орон зайн онцлог шинжүүдийг автоматаар ялган сурах чадвартай тул эмнэлгийн дүрсний ангилал болон сегментчлэлд өргөн ашиглагддаг.

Уламжлалт нейрон сүлжээнүүд нь оролтын бүх пикселийг бүрэн холбогдсон давхаргаар боловсруулдаг тул өндөр хэмжээст дүрс дээр параметрийн тоо огцом өсдөг. Харин CNN нь convolution үйлдлийг ашигласнаар параметрийн тоог бууруулж, дүрсний орон зайн мэдээллийг хадгалан боловсруулах боломжтой байдаг.

Зураг~\ref{fig:cnn_arch}-д CNN архитектурын үндсэн бүрэлдэхүүн хэсгүүдийг үзүүлэв. Convolution давхарга нь дүрсээс онцлог шинжүүдийг ялган авч, pooling давхарга нь мэдээллийн хэмжээг бууруулан илүү тогтвортой төлөөлөл үүсгэдэг. Төгсгөлд нь бүрэн холбогдсон давхарга ангиллын үр дүнг гаргадаг.

Конволюцийн үйлдлийг дараах байдлаар илэрхийлдэг.

\begin{equation}
    S(i,j) = (I * K)(i,j) = \sum_{m}\sum_{n} I(i+m,\, j+n)\, K(m,n)
    \label{eq:convolution}
\end{equation}

Энд $I$ нь оролтын дүрс, $K$ нь convolution kernel, харин гаралтын утга $S(i,j)$ нь kernel болон дүрсний локал хэсгийн үржвэрийн нийлбэрээр тодорхойлогдоно.

Ангиллын даалгаварт CNN-ийн эцсийн давхаргаас гарсан логит утгуудыг магадлал болгон хувиргахад softmax функцийг ашиглана. $C$ ангилалтай үед $c$ дугаар ангиллын магадлалыг дараах байдлаар илэрхийлнэ.

\begin{equation}
    p_c = \frac{\exp(z_c)}{\sum_{k=1}^{C}\exp(z_k)}
    \label{eq:softmax}
\end{equation}

Энд $z_c$ нь $c$ дугаар ангиллын логит утга, $p_c$ нь тухайн ангилалд хамаарах таамагласан магадлал юм. Загварыг сургах үед бодит шошго $y_c$ болон таамагласан магадлал $p_c$-ийн зөрүүг cross-entropy алдагдлаар хэмжинэ.

\begin{equation}
    \mathcal{L}_{CE} = -\sum_{c=1}^{C} y_c \log(p_c)
    \label{eq:cross_entropy}
\end{equation}

Хоёртын ангиллын үед буюу хугарал байгаа эсэхийг ялгахад binary cross-entropy алдагдлыг хэрэглэж болно.

\begin{equation}
    \mathcal{L}_{BCE} = -\left[y\log(\hat{y}) + (1-y)\log(1-\hat{y})\right]
    \label{eq:bce}
\end{equation}

Энд $y$ нь бодит шошго, $\hat{y}$ нь загварын таамагласан магадлал юм. CNN нь эмнэлгийн дүрс дэх хугарал, хавдар, эмгэг өөрчлөлт зэрэг бүтэцтэй холбоотой мэдээллийг автоматаар сурах чадвартай тул энэхүү судалгааны ангиллын модулийн онолын үндэс болж байна.

\subsection*{Шилжүүлэн суралцах арга}

\begin{figure}[H]
    \centering
    \figauto[0.60\textwidth]{Figures/transfer_2.2.png}{Transfer Learning-ийн схем (албан ёсны эх сурвалжаас)}
    \figcaption{Transfer Learning-ийн схем}{Transfer Learning-ийн схем — том өгөгдөл (ImageNet) дээр урьдчилан сургасан загварын жинг зорилтот эмнэлгийн дүрсэнд дахин ашиглаж нарийн тохируулдаг (fine-tuning) зарчим}
    \label{fig:transfer_learning}
\end{figure}

Эмнэлгийн дүрсний өгөгдөл цуглуулах болон тэмдэглэх үйл ажиллагаа ихээхэн цаг хугацаа, хүний нөөц шаарддаг. Иймээс ихэнх эмнэлгийн судалгаанд Transfer Learning буюу шилжүүлэн суралцах аргыг ашигладаг~\cite{ref30}.

Шилжүүлэн суралцах арга нь том хэмжээний өгөгдлийн сан дээр урьдчилан сургасан загварын мэдлэгийг шинэ даалгаварт ашиглах зарчимд суурилдаг. Ихэвчлэн ImageNet өгөгдлийн сан дээр сургасан жингүүдийг эмнэлгийн дүрсний ангилалд дахин ашигладаг.

Зураг~\ref{fig:transfer_learning}-д Transfer Learning-ийн ерөнхий процессыг харуулав. Эхлээд ерөнхий зориулалтын өгөгдөл дээр сурсан загварын жинг авч, дараа нь зорилтот өгөгдөл дээр нарийн тохируулга (fine-tuning) хийдэг.

Энэхүү арга нь сургалтын хугацааг бууруулахаас гадна бага хэмжээний өгөгдөлтэй нөхцөлд илүү өндөр гүйцэтгэл үзүүлэх боломж олгодог.

\subsection*{ConvNeXt архитектур}

\begin{figure}[H]
    \centering
    \figauto[0.42\textwidth]{Figures/ConvNeXt-Base_2.3.jpg}{ConvNeXt-Base архитектурын ерөнхий бүтэц}
    \figcaption{ConvNeXt-Base архитектурын ерөнхий бүтэц}{ConvNeXt-Base архитектурын ерөнхий бүтэц — depthwise convolution, LayerNorm болон GELU ашигладаг, дөрвөн шаттай иерархи бүтэцтэй орчин үеийн CNN архитектур}
    \label{fig:convnext_block}
\end{figure}

ConvNeXt нь орчин үеийн CNN архитектур бөгөөд уламжлалт convolution сүлжээний давуу талыг хадгалан Vision Transformer архитектурын зарим санааг шингээсэн байдаг~\cite{liu2022convnext}.

ConvNeXt архитектурын суурь бүтэц нь дөрвөн шаталсан үе шат (stage)-аас бүрдэнэ. Үе шат бүр орон зайн нягтаршлыг хоёр дахин бууруулж, сувгийн тоог нэмэгдүүлснээр дүрслэлийн шинжийг бүдүүвчээс нарийн руу шаталсан байдлаар (иерархи) сурдаг. Энэ нь ResNet-ийн \texttt{res2--res5} үе шаттай төстэй боловч блокийн дотоод бүтэц нь Vision Transformer-ийн зарчмыг тусгасан байдгаараа ялгаатай.

Блок бүр нь \emph{residual} холболттой буюу орцыг гаралттай нэмж дамжуулдаг ($x + \mathcal{F}(x)$). Ингэснээр гүн сүлжээнд градиент алга болох (vanishing gradient) асуудлыг бууруулж, мэдээллийг шууд урсгах боломж олгодог. Блокийн дотор эхлээд $7\times7$ хэмжээтэй \emph{depthwise convolution} үйлдлийг хэрэглэнэ; энэ нь суваг (channel) тус бүрд тусад нь convolution хийснээр энгийн convolution-той харьцуулахад параметр, тооцооллыг эрс багасгаж, том хүлээн авах талбай (receptive field) олж авах боломж бүрдүүлдэг. Үүний дараа сувгийн тоог түр нэмэгдүүлээд дахин буулгадаг \emph{урвуу bottleneck} (inverted bottleneck) бүтэцтэй $1\times1$ convolution-ууд дамжина.

Хэвийн болголтод (normalization) уламжлалт Batch Normalization-ийн оронд \emph{Layer Normalization} (LayerNorm)-ийг ашигладаг бөгөөд энэ нь batch-ийн хэмжээнээс хамаардаггүй тул жижиг batch-тай сургалтад илүү тогтвортой. Идэвхжүүлэлтийн функцэд ReLU-гийн оронд гөлгөр шинжтэй \emph{GELU} (Gaussian Error Linear Unit)-ийг хэрэглэх ба блок бүрт цөөн тооны идэвхжүүлэлт болон normalization давхарга байрлуулснаар Transformer блокийн загварт ойртуулсан. GELU-г дараах байдлаар ойролцоолон илэрхийлнэ.

\begin{equation}
    \mathrm{GELU}(x) = x\,\Phi(x) \approx 0.5\,x\left(1 + \tanh\!\left[\sqrt{2/\pi}\,\bigl(x + 0.044715\,x^{3}\bigr)\right]\right)
    \label{eq:gelu}
\end{equation}

Энд $\Phi(x)$ нь стандарт нормаль тархалтын хуримтлагдсан функц юм. Эдгээр шийдлүүдийн нийлбэр нь ConvNeXt-ийг convolution сүлжээний орон зайн inductive bias-ийг хадгалангаа Vision Transformer-ийн нарийвчлалд хүргэх боломжийг олгодог. Тус архитектур нь ImageNet зэрэг том хэмжээний өгөгдөл дээр өндөр үр дүн үзүүлсэн бөгөөд эмнэлгийн дүрсний ангиллын олон судалгаанд ашиглагдаж байна.

Энэхүү судалгаанд ConvNeXt төрлийн загваруудыг харьцуулан туршиж, эцсийн ангиллын модульд ConvNeXt-Base хувилбарыг ашигласан. Том хэмжээний ImageNet-22K урьдчилсан сургалт болон Base архитектурын илүү өндөр төлөөлөх чадвар нь рентген зураг дахь нарийн хугарлын шинжийг ялгахад тохиромжтой гэж үзсэн.

\subsection*{Дүрсний сегментчлэл ба U-Net архитектур}

\begin{figure}[H]
    \centering
    \figauto[0.62\textwidth]{Figures/u_net_2.4.png}{U-Net архитектур (албан ёсны эх сурвалжаас)}
    \figcaption{U-Net архитектур}{U-Net архитектур — кодлогч (encoder) ба тайлагч (decoder) хэсгээс бүрдэх бөгөөд skip connection-оор доод түвшний орон зайн мэдээллийг хадгалж жижиг хугарлын бүсийг нарийвчлан сэргээдэг сегментчлэлийн сүлжээ}
    \label{fig:unet}
\end{figure}

Дүрсний сегментчлэл нь дүрсний пиксел бүрийг тодорхой ангилалд хуваарилах процесс юм. Эмнэлгийн дүрс боловсруулалтад сегментчлэл нь эмгэг өөрчлөлтийн байршил болон хэлбэрийг нарийвчлан тодорхойлоход ашиглагддаг.

Сегментчлэлийн чанарыг үнэлэхэд Dice coefficient болон Intersection over Union (IoU) хэмжүүрүүд өргөн хэрэглэгддэг~\cite{ref13}.

\begin{equation}
    \mathrm{Dice}(A, B) = \frac{2\,|A \cap B|}{|A| + |B|}
    \label{eq:dice}
\end{equation}

\begin{equation}
    \mathrm{IoU}(A, B) = \frac{|A \cap B|}{|A \cup B|}
    \label{eq:iou}
\end{equation}

Энд $A$ нь бодит маск, $B$ нь таамагласан маск. Dice нь бодит болон таамагласан маскуудын давхцлыг хэмждэг бол IoU нь тэдгээрийн огтлолцол болон нийлбэр мужийн харьцааг илэрхийлдэг.

Сургалтын үед сегментчлэлийн загварын гаралт нь тасралтгүй магадлалын зураг хэлбэртэй байдаг тул Dice loss-ийг пикселийн түвшний магадлал дээр дараах байдлаар тооцно.

\begin{equation}
    \mathcal{L}_{Dice} =
    1 - \frac{2\sum_{i=1}^{N} p_i g_i + \varepsilon}
    {\sum_{i=1}^{N} p_i + \sum_{i=1}^{N} g_i + \varepsilon}
    \label{eq:dice_loss}
\end{equation}

Энд $p_i$ нь $i$ дугаар пикселийн таамагласан магадлал, $g_i$ нь бодит маскны утга, $N$ нь нийт пикселийн тоо, $\varepsilon$ нь хуваарь тэг болохоос сэргийлэх бага тогтмол юм. Жижиг хугарлын бүс нийт дүрстэй харьцуулахад маш бага талбай эзэлдэг тул Dice loss нь ангийн тэнцвэргүй байдлын нөлөөг бууруулахад тохиромжтой.

U-Net нь эмнэлгийн дүрсний сегментчлэлд зориулан боловсруулсан архитектур бөгөөд кодлогч (encoder) болон тайлагч (decoder) хэсгээс бүрддэг~\cite{ronneberger2015unet}.

Skip connection ашигласнаар дүрсний нарийн мэдээллийг хадгалж үлддэг бөгөөд энэ нь жижиг хэмжээний эмгэг өөрчлөлтийг ялгахад чухал үүрэгтэй.

\subsection*{Segment Anything Model ба LoRA}

\begin{figure}[H]
    \centering
    \figauto[0.62\textwidth]{Figures/SAM_REAL_LAST_2.5.jpg}{SAM архитектур (албан ёсны эх сурвалжаас)}
    \figcaption{SAM архитектур}{SAM архитектур — image encoder, prompt encoder болон mask decoder-оос бүрдэх суурь загвар; энэ судалгаанд YOLO хайрцаг тэмдэглэгээг prompt болгон ашиглаж хугарлын псевдо маск үүсгэхэд хэрэглэв}
    \label{fig:sam}
\end{figure}

Segment Anything Model (SAM) нь Meta AI байгууллагаас боловсруулсан foundation model бөгөөд олон төрлийн объектыг сегментчлэх чадвартай~\cite{ref32}.

SAM нь image encoder, prompt encoder болон mask decoder гэсэн үндсэн хэсгүүдээс бүрдэнэ. Урьдчилан сургасан том хэмжээний мэдлэгийн ачаар бага хэмжээний тэмдэглэгээтэй өгөгдөл дээр ч ерөнхий сегментчлэлийн чадвар үзүүлэх боломжтой.

Энэхүү судалгаанд SAM-ийг эцсийн сегментчлэлийн үндсэн загвар болгон бус, харин YOLO хайрцаг тэмдэглэгээг prompt болгон ашиглаж хугарлын бүсийн псевдо маск үүсгэхэд хэрэглэсэн. Үүссэн псевдо маскуудыг дараа нь ResNet50 U-Net сегментчлэгчийн сургалтын өгөгдлийг өргөтгөхөд ашиглав.

\begin{figure}[H]
    \centering
    \figauto[0.42\textwidth]{Figures/LORA_LAST_2.6.png}{LoRA архитектур (албан ёсны эх сурвалжаас)}
    \figcaption{LoRA архитектур}{LoRA архитектур — үндсэн жинг хөлдөөж зөвхөн бага зэрэглэлийн матрицуудыг ($A$, $B$) сургаснаар домэйнд тохируулах санах ой, тооцооллын өртгийг мэдэгдэхүйц бууруулдаг}
    \label{fig:lora}
\end{figure}

SAM зэрэг том загваруудыг бүрэн дахин сургах нь өндөр нөөц шаарддаг тул домэйнд тохируулах үед Low-Rank Adaptation (LoRA) зэрэг параметрийн үр ашигтай аргыг ашиглах боломжтой~\cite{hu2022lora}. Энэ судалгааны эхний олон даалгаварт туршилтад SAM-LoRA хандлагыг авч үзсэн боловч эцсийн pipeline-д сегментчлэлийн тогтвортой байдлыг нэмэгдүүлэх зорилгоор тусдаа ResNet50 U-Net загварт шилжсэн.

LoRA нь үндсэн жинг өөрчлөхгүйгээр бага хэмжээний параметр сургах замаар шинэ даалгаварт тохируулах боломж олгодог. Жинг шинэчлэх томьёог дараах байдлаар илэрхийлнэ.

\begin{equation}
    W = W_0 + \Delta W = W_0 + \frac{\alpha}{r}\, B A,
    \qquad B \in \mathbb{R}^{d \times r},\; A \in \mathbb{R}^{r \times k},\; r \ll \min(d, k)
    \label{eq:lora}
\end{equation}

Энд $W_0$ нь хөлдөөсөн үндсэн жин, $A$ ба $B$ нь сургагдах бага зэрэглэлийн (low-rank) матрицууд, $r$ нь зэрэглэл, $\alpha$ нь масштабын коэффициент юм.

LoRA нь бүх жинг дахин сургахын оронд зөвхөн бага зэрэглэлийн шинэчлэлт ($BA$)-ийг сургадаг тул сургагдах параметрийн тоо, санах ойн зарцуулалт болон тооцооллын өртгийг мэдэгдэхүйц бууруулдаг. Үндсэн жин $W_0$ хөлдөөстэй хэвээр үлддэг тул урьдчилан сурсан мэдлэг хадгалагдаж, домэйнд тохируулах сургалт нь цөөн өгөгдөл болон хязгаарлагдмал тооцооллын нөөц бүхий нөхцөлд ч тогтвортой явагдах боломжтой болдог.

\subsection*{Загварын тайлбарлах чадвар}

\begin{figure}[H]
    \centering
    \figauto[0.66\textwidth]{Figures/Grad_cam_2.7.jpg}{Grad-CAM-ийн ажиллах урсгал (албан ёсны эх сурвалжаас)}
    \figcaption{Grad-CAM-ийн ажиллах урсгал}{Grad-CAM-ийн ажиллах урсгал — градиентэд суурилан загвар шийдвэр гаргахдаа дүрсний аль хэсэгт төвлөрснийг heatmap хэлбэрээр дүрсэлдэг}
    \label{fig:gradcam}
\end{figure}

\begin{figure}[H]
    \centering
    \figauto[0.56\textwidth]{Figures/ShAP_2.8.png}{SHAP framework (албан ёсны эх сурвалжаас)}
    \figcaption{SHAP framework}{SHAP framework — тоглоомын онолын Shapley утгад үндэслэн оролтын шинж бүр таамаглалд хэр нөлөөлснийг тоон байдлаар тооцож тайлбарладаг арга}
    \label{fig:shap}
\end{figure}

Эмнэлгийн салбарт зөвхөн өндөр нарийвчлалтай байх нь хангалтгүй бөгөөд загварын шийдвэрийг тайлбарлах боломжтой байх шаардлагатай.

Grad-CAM нь CNN загварын анхаарч буй хэсгийг heatmap хэлбэрээр дүрслэх арга юм~\cite{ref37}.

Тухайн $c$ ангиллын хувьд сүүлийн convolution давхаргын $k$ дугаар feature map-ийн ач холбогдлын жинг дараах байдлаар тооцно.

\begin{equation}
    \alpha_k^c = \frac{1}{Z}\sum_i\sum_j
    \frac{\partial y^c}{\partial A_{ij}^{k}}
    \label{eq:gradcam_weight}
\end{equation}

Энд $A_{ij}^{k}$ нь $k$ дугаар feature map-ийн $(i,j)$ байрлал дахь идэвхжил, $y^c$ нь $c$ ангиллын score, $Z$ нь feature map-ийн нийт байрлалын тоо юм. Дараа нь class activation map-ийг дараах байдлаар гаргана.

\begin{equation}
    L_{\mathrm{Grad\text{-}CAM}}^{c} =
    \mathrm{ReLU}\left(\sum_k \alpha_k^c A^k\right)
    \label{eq:gradcam_map}
\end{equation}

Энэхүү арга нь загвар шийдвэр гаргахдаа дүрсний аль хэсгийг ашигласныг харуулах боломж олгодог.

Мөн SHAP арга нь тухайн таамаглалд оролтын шинжүүд ямар хэмжээний нөлөө үзүүлснийг тооцоолдог~\cite{ref36}.

SHAP нь тоглоомын онолын Shapley утгад суурилдаг бөгөөд $i$ дугаар шинжийн хувь нэмрийг дараах томьёогоор тодорхойлно.

\begin{equation}
    \phi_i =
    \sum_{S \subseteq F \setminus \{i\}}
    \frac{|S|!\,(|F|-|S|-1)!}{|F|!}
    \left[f(S \cup \{i\}) - f(S)\right]
    \label{eq:shap}
\end{equation}

Энд $F$ нь бүх шинжийн олонлог, $S$ нь $i$ шинжийг агуулаагүй дэд олонлог, $f(S)$ нь тухайн шинжүүдийг ашигласан үеийн загварын гаралт юм. SHAP нь загварын шийдвэр гаргалтыг илүү ойлгомжтой тайлбарлах боломж олгодог бөгөөд энэхүү судалгаанд эдгэрэлтийн үнэлгээний үр дүнг шинжлэхэд ашигласан.

\subsection*{Үнэлгээний хэмжүүрүүдийн онолын томьёо}

Ангиллын загварын үр дүнг үнэлэхэд confusion matrix-д суурилсан хэмжүүрүүдийг ашиглана. Энд $TP$ нь зөв эерэг, $TN$ нь зөв сөрөг, $FP$ нь буруу эерэг, $FN$ нь буруу сөрөг тохиолдлын тоог илэрхийлнэ.

\begin{equation}
    \mathrm{Accuracy} = \frac{TP + TN}{TP + TN + FP + FN}
    \label{eq:accuracy}
\end{equation}

\begin{equation}
    \mathrm{Precision} = \frac{TP}{TP + FP},
    \qquad
    \mathrm{Recall} = \frac{TP}{TP + FN}
    \label{eq:precision_recall}
\end{equation}

\begin{equation}
    F_1 = 2 \cdot
    \frac{\mathrm{Precision}\cdot\mathrm{Recall}}
    {\mathrm{Precision}+\mathrm{Recall}}
    \label{eq:f1}
\end{equation}

Accuracy нь нийт зөв таамаглалын хувийг, Precision нь хугарал гэж таамагласан тохиолдлуудын хэд нь үнэн зөв болохыг, Recall нь бодит хугаралтай тохиолдлын хэдийг загвар илрүүлснийг тус тус илэрхийлнэ. Харин $F_1$ нь Precision болон Recall-ийн тэнцвэрийг нэгтгэн харуулдаг тул ангийн тэнцвэргүй өгөгдөлтэй үед илүү мэдээлэлтэй хэмжүүр болдог.


%==============================================================
\section{Ижил төстэй судалгаанууд}
%==============================================================

\subsection*{Гадаадын судалгаанууд}

Ясны хугарлыг хиймэл оюуны аргаар илрүүлэх чиглэлээр олон судалгаа хийгдсэн. Lindsey нарын боловсруулсан загвар нь насанд хүрэгчдийн яс-булчингийн рентген зургууд дээр 16 анатомийн бүсэд хугарал илрүүлэх даалгаврыг гүйцэтгэж, AUC 0.948--0.982 үзүүлэлтэд хүрсэн байна~\cite{ref08}. Энэ ажил нь гүн сургалтын загвар радиологичийн шийдвэрийг дэмжих боломжтойг харуулсан боловч зөвхөн хугарал илрүүлэх асуудалд төвлөрсөн бөгөөд хугарлын байршил, хугарлын зэрэг болон эдгэрэлтийн үнэлгээг хамраагүй.

Yacoub нарын судалгаа нарийн болон бүдэг хугарлыг илрүүлэхэд гүн сургалтын загвар ашиглах боломжийг авч үзсэн бөгөөд уламжлалт рентген уншилтад алдагдах магадлалтай хугарлын төрлүүд дээр AI-ийн дэмжлэг хэрэгтэйг онцолсон~\cite{ref10}. Kitamura нар шагайн хугарлыг илрүүлэх CNN ensemble загварыг бага хэмжээний өгөгдөл дээр сургаж, олон өнцгийн рентген зураг ашиглах нь гүйцэтгэлд нөлөөлдгийг харуулсан~\cite{ref14}. Tanzi нар proximal femur хугарлыг шаталсан гүн сургалтын аргаар ангилах судалгаа хийж, зөвхөн хугарал байгаа эсэхээс гадна хугарлын төрлийг ялгах асуудлыг авч үзсэн~\cite{ref15}.

Мөн Gu болон Zheng~\cite{ref13} нар ясны хугарлын оношилгоонд хиймэл оюун ухааныг ашигласан 337 судалгаанаас 40 өгүүллийг сонгон системтэй тойм судалгаа хийжээ. Тэд хугарал илрүүлэх, ангилах, сегментлэх зэрэг даалгавруудыг тусад нь ангилан судалсан бөгөөд олон даалгаварт загвар болон загварын тайлбарлах чадварын чиглэлд судалгаа харьцангуй хомс байгааг онцолсон байна. Түүнчлэн ихэнх судалгаа зөвхөн нэг төрлийн даалгаварт төвлөрсөн бөгөөд клиникийн хэрэглээнд шаардлагатай нэмэлт мэдээллийг хангалттай авч үзээгүй болохыг дүгнэжээ.

Сүүлийн жилүүдэд зөвхөн загварын архитектур төдийгүй өгөгдлийн сангийн чанар, тэмдэглэгээний төрөл чухал асуудал болж байна. Rajpurkar нарын MURA өгөгдлийн сан нь яс-булчингийн рентген зургаар abnormality detection хийх томоохон эх сурвалж болсон~\cite{ref39}. Abedeen нарын FracAtlas өгөгдөл нь хугарлын ангилал, байршил болон сегментчлэлийн тэмдэглэгээг хамтад нь өгснөөр хугарал илрүүлэхээс гадна хугарлын бүсийг пикселийн түвшинд судлах боломж олгосон~\cite{ref40}. Nagy нарын GRAZPEDWRI-DX өгөгдөл нь хүүхдийн бугуйн гэмтлийн рентген зураг, YOLO тэмдэглэгээ болон мета-өгөгдөл агуулдаг тул хүүхдийн хугарлын автомат оношилгооны судалгаанд ач холбогдолтой~\cite{ref41}.

\subsection*{Монгол Улс дахь судалгаанууд}

Монгол Улсад ясны хугарлыг рентген зураг болон гүн сургалтын аргаар оношлох чиглэлийн судалгаа одоогоор хязгаарлагдмал байна. Ganmaa нар~\cite{ref20} Улаанбаатар хотын сурагчдыг хамруулсан гурван жилийн хугацаатай судалгаагаар хүүхдийн хугарлын тархалт болон эрсдэлийн талаар үнэлгээ хийсэн бөгөөд Монгол хүүхдийн дундах хугарлын эпидемиологийн мэдээллийг бүрдүүлсэн. Гэвч энэхүү судалгаанд дүрс боловсруулалт болон хиймэл оюуны арга ашиглаагүй болно.

Мөн Jaalkhorol нар~\cite{ref21} Монгол хүн амд зориулсан FRAX загварыг нутагшуулж, ясны хугарлын эрсдэлийг үнэлэх судалгаа хийсэн. Уг ажил нь үндэсний түвшний эрсдэлийн үнэлгээнд ач холбогдолтой боловч рентген зурагт суурилсан автомат оношилгоо болон гүн сургалтын аргуудыг хамраагүй байна.

\subsection*{Гадаад болон Монголын судалгааны харьцуулалт}

Ижил төстэй судалгаануудыг нэгтгэн дүгнэхэд гадаадын судалгаанд хугарал илрүүлэх, ангилах чиглэлийн гүн сургалтын аргууд эрчимтэй хөгжиж байгаа боловч хугарлыг илрүүлэх, сегментлэх, хугарлын зэргийг тодорхойлох болон эдгэрэлтийн үнэлгээг нэгтгэн авч үзсэн судалгаа харьцангуй цөөн байна. Харин Монгол Улсад ясны хугарлын эпидемиологи болон эрсдэлийн үнэлгээтэй холбоотой судалгаа хийгдсэн боловч рентген зурагт суурилсан хиймэл оюуны оношилгооны чиглэлд судалгаа хомс хэвээр байна. Гол шалгуурын дагуу нэгтгэсэн харьцуулалтыг Хүснэгт~\ref{tab:lit_compare}-д үзүүлэв.

\begin{table}[H]
    \centering
    \caption{Ижил төстэй судалгаануудын харьцуулалт}
    \label{tab:lit_compare}
    \small
    \begin{tabular}{p{2.7cm} p{2.6cm} p{2.7cm} p{2.7cm} p{3.0cm}}
    \toprule
    \textbf{Судалгаа} & \textbf{Өгөгдөл} & \textbf{Даалгавар} & \textbf{Арга / Загвар} & \textbf{Хязгаарлалт} \\
    \midrule
    Lindsey нар~\cite{ref08} & Яс-булчингийн рентген зураг & Хугарал илрүүлэх & Deep CNN & Сегментчлэл, хугарлын зэрэг, эдгэрэлт хамраагүй \\[4pt]
    Yacoub нар~\cite{ref10} & Нарийн хугарлын рентген зураг & Далд/бүдэг хугарал илрүүлэх & Deep learning & Нэг даалгаварт төвлөрсөн \\[4pt]
    Kitamura нар~\cite{ref14} & Шагайн рентген зураг & Хугарал илрүүлэх & CNN ensemble & Өгөгдлийн хэмжээ бага, нэг анатомийн бүс \\[4pt]
    Tanzi нар~\cite{ref15} & Proximal femur X-ray & Хугарлын төрөл ангилах & Multistage deep learning & Сегментчлэл, эдгэрэлт ороогүй \\[4pt]
    Rajpurkar нар~\cite{ref39} & MURA & Abnormality detection & CNN суурь загварууд & Abnormal шошго нь зөвхөн хугарал биш \\[4pt]
    Abedeen нар~\cite{ref40} & FracAtlas & Ангилал, байршил, сегментчлэл & Dataset benchmark & Хугарлын зэргийн бүрэн клиникийн шошго хязгаарлагдмал \\[4pt]
    Nagy нар~\cite{ref41} & GRAZPEDWRI-DX & Хүүхдийн бугуйн гэмтэл & YOLO тэмдэглэгээтэй өгөгдөл & Нэг анатомийн бүс, эдгэрэлтийн шошго байхгүй \\[4pt]
    Энэхүү судалгаа & Олон эх сурвалжийн рентген зураг & Илрүүлэх, сегментлэх, хугарлын зэрэг, эдгэрэлт & ConvNeXt-Base, ResNet50 U-Net, морфологийн шинж & Эдгэрэлт синтетик өгөгдөлд тулгуурласан \\
    \bottomrule
    \end{tabular}
\end{table}

Хүснэгт~\ref{tab:lit_compare}-ээс харахад гадаадын судалгаанд хугарал илрүүлэх болон тодорхой төрлийн хугарлыг ангилах чиглэлийн ажил харьцангуй хөгжсөн боловч хугарлын байршил, хугарлын зэрэг, эдгэрэлтийн таамаглал болон тайлбарлах чадварыг нэг pipeline-д нэгтгэн авч үзсэн ажил хязгаарлагдмал байна. Харин Монгол Улсад рентген зурагт суурилсан хиймэл оюуны оношилгооны судалгаа бараг байхгүй байгаа нь энэ чиглэлд цаашид судалгаа хийх шаардлагатайг харуулж байна.

%==============================================================
\section{Судалгааны цоорхой}
%==============================================================

Холбогдох судалгаануудын дүн шинжилгээнээс харахад ясны хугарлыг хиймэл оюуны аргаар оношлох чиглэлд мэдэгдэхүйц ахиц гарсан хэдий ч хэд хэдэн асуудал бүрэн шийдэгдээгүй хэвээр байна.

Нэгдүгээрт, ижил төстэй судалгааны ихэнх нь хугарал илрүүлэх, ангилах эсвэл сегментлэх зэрэг нэг төрлийн даалгаварт төвлөрсөн байдаг бөгөөд эдгээрийг нэгтгэн авч үзсэн судалгаа харьцангуй ховор байна~\cite{ref13}. Ялангуяа хугарал илрүүлэх, сегментлэх, хугарлын зэргийг тодорхойлох болон эдгэрэлтийн үнэлгээг нэг хүрээнд авч үзсэн ажил бараг байхгүй байна.

Хоёрдугаарт, эдгээр даалгавруудыг нэг хуваалцсан кодлогчтой олон даалгаварт (multi-task)~\cite{ref42} загварт нэгтгэх оролдлогууд хийгдэж байгаа боловч даалгавруудын хооронд үүсэх task interference болон түүнийг бууруулах модульчлагдсан (modular) хандлагыг ясны хугарлын домэйнд системтэйгээр харьцуулан судалсан ажил ховор байна. Мөн Segment Anything Model (SAM) зэрэг суурь загварыг эмнэлгийн дүрсэнд шууд эцсийн сегментлэгч бус, харин сургалтын псевдо маск үүсгэгч болгон ашиглах болон LoRA зэрэг параметрийн үр ашигтай тохируулгын аргыг ясны хугарлын өгөгдөлд хэрэглэсэн судалгаа хязгаарлагдмал хэвээр байна~\cite{ref15}.

Гуравдугаарт, хугарлын зэргийг автоматаар үнэлэх нэгдсэн аргачлал бүрэн тогтоогдоогүй бөгөөд ихэнх тохиолдолд эмчийн субъектив үнэлгээнд тулгуурладаг байна~\cite{ref14}. Үүний улмаас судалгааны үр дүнг хооронд нь харьцуулах болон стандартчилахад хүндрэл үүсдэг.

Дөрөвдүгээрт, ясны хугарлын оношилгоог эдгэрэлтийн үнэлгээ, таамаглалтай уялдуулан судалсан ажил харьцангуй цөөн байна. Ихэнх судалгаа зөвхөн оношилгооны шатанд төвлөрсөн байдаг бөгөөд эмчилгээний дараах явцыг үнэлэх асуудлыг хангалттай авч үзээгүй байна.

Монгол Улсын хувьд ясны хугарлын тархалт болон эрсдэлийн талаарх судалгаанууд хийгдсэн боловч рентген зурагт суурилсан хиймэл оюуны оношилгооны чиглэлээр хийгдсэн судалгаа бараг байхгүй байна~\cite{ref20,ref21}. Иймээс энэхүү чиглэлээр судалгаа хийх нь олон улсын судалгааны цоорхойг нөхөхөөс гадна дотоодын нөхцөлд хиймэл оюуны хэрэглээг өргөжүүлэх ач холбогдолтой юм.

Иймд энэхүү судалгаанд ясны хугарлыг илрүүлэх, сегментлэх, хугарлын зэргийг үнэлэх болон эдгэрэлтийн таамаглалтай холбоотой асуудлуудыг нэг хүрээнд авч үзэж, ижил төстэй судалгаануудад тодорхойлогдсон дээрх цоорхойг нөхөхийг зорьсон болно.
